\documentclass[a4paper,10pt]{article}
\usepackage[left=1in, right=1in, top=1in, bottom=1in]{geometry}
\usepackage{enumitem, graphicx, hyperref}
\usepackage{xcolor}
\usepackage{eso-pic,transparent}

\AddToShipoutPicture{
    \AtPageLowerLeft{
        \raisebox{0.25\paperheight}{ % Adjust this value to change the vertical position
            \hspace{0.1\paperwidth} % Adjust this value to change the horizontal position
            
\includegraphics[width=0.80\paperwidth]{logo 素材.jpg} % Replace with your background image file
        }
    }
}

\begin{document}

\pagestyle{empty}

\begin{center}
    \begin{minipage}{0.75\textwidth}
        {\LARGE \textbf{Curriculum Vitae}}\\ %Date of Birth: 2000.02.04\\  % I don't want to show DOB, which shows that I am too old
        {\large Name: Liang Yue \\ Email: yueliang200024@gmail.com \\ Phone: 17801014228 }\\
        {\large GitHub: \href{https://github.com/Philip-2000}{https://github.com/Philip-2000} }
    \end{minipage}
    \begin{minipage}{0.20\textwidth}
        \begin{flushright}
            \includegraphics[width=2.0cm]{IMG_2160.jpg} % Replace "yourphoto.jpg" with your actual photo file name
        \end{flushright}
    \end{minipage}
\end{center}

\hrule

\section*{Education}

\noindent \textbf{Tsinghua University, Department of Computer Science and Technology}\\
\textit{Master of Science in Computer Science} \hfill {Expected Graduation: June 2025}\\
Advised by Professor Song-Hai Zhang, conducting research on "Digital 3D Indoor Scene Construction Algorithms" in the Graphics and Geometric Computing Laboratory.

\vspace{0.5em}
\noindent \textbf{Peking University, School of Information Science and Technology}\\
\textit{Bachelor of Science in Computer Science, Double Major in Physics} \hfill {Graduation: June 2022}\\
Relevant Courses: Mathematical Analysis, Computer Graphics, Artificial Intelligence, Operating Systems, etc.\\

\hrule

\section*{Publications}
\begin{itemize}[leftmargin=*]
    \item Liang Yue, Hao Tan, Junkai Huang, Shao-Kui Zhang. \textbf{A Survey on Digital 3D Indoor Scene Construction}, Journal of Image and Graphics, 2024.
    \item Shao-Kui Zhang, Junkai Huang, Liang Yue, Jia-Tong Zhang, Jia-Hong Liu, Yu-Kun Lai, Song-Hai Zhang. \textbf{SceneExpander: Real-Time Scene Synthesis for Interactive Floor Plan Editing}, ACM Multimedia, 2024. \href{https://dl.acm.org/doi/10.1145/3664647.3680798}{doi:10.1145/3664647.3680798}
    \item Liang Yue, Shao-kui Zhang, Lin Yuan, Yi-tao Chen, Zi-rui Zhou, Song-Hai Zhang. \textbf{Synthesizing 3D Scenes via Diffusion Model that Incorporates Indoor Scene Characteristics}, ACM Multimedia, 2025. \href{https://dl.acm.org/doi/10.1145/3746027.3754832}{doi:10.1145/3746027.3754832}
    \item Shao-kui Zhang, Liang Yue, Han-xi Zhu, Hao-ran Zhang, Song-Hai Zhang. \textbf{SceneCluster: Interactive Scene Synthesis by Clustering Groups of Furniture Objects}, IEEE TVCG, 2025. \href{https://ieeexplore.ieee.org/document/11260934}{doi:10.1109/TVCG.2025.3634829}
\end{itemize}

\hrule

\section*{Internships}
\begin{itemize}[leftmargin=*]
    \item Tsinghua University, College of Artificial Intelligence — Research Intern (Advisor: Miao Liu). 2025-06 -- Present.
    \item Kunlun Skyworks — Project: World Model. 2025-07 -- 2025-09. (Finetuning navigation world model; evaluation on familiar / half-familiar / unfamiliar scenes.)
    \item RayNeo Smart Glasses (TCL) — Project: Personal Egocentric Video Assistant. 2025-10 -- Present. (Memory construction and retrieval for QA.)
\end{itemize}

\hrule

\section*{Projects}

\vspace{0.5em}
\textbf{Finetuning Multimodal Large Language Model with Egocentric Video Data}\\
Enhancing the model's capability in understanding and reasoning about first-person visual information through egocentric video data. (2026-01 -- till now)\\

\vspace{0.5em}
\noindent \textbf{Personal Egocentric Video Assistant}\\
Developing a personal egocentric video assistant system with a focus on memory construction and retrieval for question answering. (2025-10 -- till now)\\

\vspace{0.5em}
\noindent \textbf{Finetuning Navigation World Model with New Datasets}\\
Finetuning the Navigation World Model using the SEKAI navigation dataset and evaluating its performance across familiar, half-familiar, and unfamiliar scenes. (2025-07 -- 2025-09)\\

\vspace{0.5em}
\noindent \textbf{Field-based Indoor Scene Optimization Methods}\\
Integrated features like physical constraints and layout rationality into a guiding "field" for object layout optimization. (2024-11 -- 2025-04)\\

\vspace{0.5em}
\noindent \textbf{Semantics for Indoor Scene Layout and Dataset Augmentation}\\
Digital indoor scenes are complex data structures; we model semantics such as "object relationships" and "functional groups" to expand indoor scene datasets for training deep generative models. (2024-07 -- 2024-11)\\

\hrule

\section*{Research Interests}
\begin{itemize}[leftmargin=*]
    \item \textbf{Augmented Reality.} Next generation of Human Computer Interaction.
    \item \textbf{3D Computer Vision.} Mixture of Virtual and Real World.
    \item \textbf{Computer Graphics.} Digital Content Creation and Interaction.
\end{itemize}


\end{document}
